% GENERATED by docs/latex/tools/gen_cli_manual_tables.sh, i.e.
% `CLI_MANUAL_EMIT=<dir> hacs CliManualCoverageValidation`, from the
% registries of the build it runs in.  Do not edit by hand: regenerate when a
% verb, type, setting, validation or command-line token changes.
\octype{Fp}{an element of F\_p (plan Sec. 3.4)}{\index{Fp@\texttt{Fp} (type)}}
\begin{membertable}
\texttt{.p} & structural & the modulus \\
\texttt{.value} & structural & the canonical residue in [0, p) \\
\end{membertable}
\octype{Poly}{a polynomial over F\_ell; prints HIGH -\textgreater{} LOW degree, while coeffs / hash use the machine order LOW -\textgreater{} HIGH (plan Sec. 2 C5)}{\index{Poly@\texttt{Poly} (type)}}
\begin{membertable}
\texttt{.ell} & structural & the modulus \\
\texttt{.deg} & structural & the degree (-1 for the zero polynomial) \\
\texttt{.monic} & structural & whether the leading coefficient is 1 \\
\end{membertable}
\octype{Matrix}{a dense matrix over F\_p (FlatMatrixT\textless{}uint32\_t\textgreater{})}{\index{Matrix@\texttt{Matrix} (type)}}
\begin{membertable}
\texttt{.rows} & structural & the row count \\
\texttt{.cols} & structural & the column count \\
\texttt{.p} & structural & the modulus \\
\texttt{.entry(row,\allowbreak{} col)} & structural & the (row, col) entry as a canonical residue \\
\end{membertable}
\octype{Model}{the SL(2,Z) orbifold model: its triangulation and the exact tessellation tables}{\index{Model@\texttt{Model} (type)}}
\begin{membertable}
\texttt{.name} & structural & the descriptor name \\
\texttt{.dims} & structural & the simplex counts by dimension \\
\texttt{.vertices} & structural & the vertex count \\
\end{membertable}
\octype{Complex}{a simplicial complex (the presets of SimplicialCohomology.h and complex(...))}{\index{Complex@\texttt{Complex} (type)}}
\begin{membertable}
\texttt{.name} & structural & the complex's name \\
\texttt{.dim} & structural & the top dimension \\
\texttt{.euler} & structural & the Euler characteristic \\
\texttt{.dims} & structural & the simplex counts by dimension \\
\texttt{.simplices(q)} & structural & the q-simplices, each as a List of vertex indices \\
\end{membertable}
\octype{LocalSystem}{Sym\textasciicircum{}n(F\_ell\textasciicircum{}2) on a model, with its rho\_n(g\_r) relations (lazy)}{\index{LocalSystem@\texttt{LocalSystem} (type)}}
\begin{membertable}
\texttt{.n} & structural & the symmetric power \\
\texttt{.p} & structural & the coefficient prime ell \\
\texttt{.d\_V} & structural & dim Sym\textasciicircum{}n = n + 1 \\
\texttt{.coefficient\_ring} & structural & the ring's canonical text: F\_ell, or Z/ell\textasciicircum{}\textless{}k\textgreater{} ([E] sec:ring (a)) \\
\texttt{.relations} & computed & the number of local-system relations \\
\end{membertable}
\octype{Cohomology}{H\textasciicircum{}q\_eq(K, R; L) with its C / Z / B / H bases and projectors (lazy); over Z/ell\textasciicircum{}k (coefficient\_ring, ell \textgreater{}= 5) its lengths, invariant factors and K-mode}{\index{Cohomology@\texttt{Cohomology} (type)}}
\begin{membertable}
\texttt{.computed} & structural & whether compute\_cohomology() has run \\
\texttt{.dims} & computed & dim H\textasciicircum{}q for every degree q \\
\texttt{.ell} & computed & the coefficient prime \\
\texttt{.d\_V} & computed & the representation dimension \\
\texttt{.q\_max} & computed & the top degree of THIS complex (derived) \\
\texttt{.truncated} & computed & false on a finite model \\
\texttt{.q\_max\_source} & computed & 'derived' unless truncated (E9) \\
\texttt{.degrees} & computed & [0 .. q\_max], derived \\
\texttt{.coefficient\_ring} & structural & the ring's canonical text: F\_ell, or Z/ell\textasciicircum{}\textless{}k\textgreater{} ([E] sec:ring (a)) \\
\texttt{.lengths} & computed & the composition length of H\textasciicircum{}q over Z/ell\textasciicircum{}k, q = 0, 1, 2 (lengths\_H\textasciicircum{}q) \\
\texttt{.ladic\_type\_certified} & computed & whether the K-mode certifies the Z\_(ell)-type at K = k ([E] sec:ring (d)) \\
\texttt{.precision} & computed & on a digit handle (ell\textasciicircum{}k \textgreater{}= 2\textasciicircum{}64): the certified precision bound -- \{precision\_requested, precision\_working, e\_max, word\_level, precision\_certified, digit\_width, digits\_L\} \\
\texttt{.bockstein\_rank} & computed & bockstein\_rank\textasciicircum{}0: the rank of beta\textasciicircum{}0 on H\textasciicircum{}0(F\_ell), from the Z/ell\textasciicircum{}2 and F\_ell objects ([E] sec:ring (d), (g)) \\
\texttt{.higher\_torsion\_count} & computed & dim H\textasciicircum{}0(F\_ell) - bockstein\_rank\textasciicircum{}0 ([E] sec:ring (d)) \\
\texttt{.C(q)} & computed & a basis of C\textasciicircum{}q\_eq, one N\_q x d\_V cochain per element \\
\texttt{.Z(q)} & computed & a basis of Z\textasciicircum{}q\_eq \\
\texttt{.B(q)} & computed & a basis of B\textasciicircum{}q\_eq \\
\texttt{.H(q)} & computed & a basis of H\textasciicircum{}q\_eq \\
\texttt{.P(q)} & computed & the projector C\textasciicircum{}q\_flat -\textgreater{} H\textasciicircum{}q coordinates (legacy shape conventions) \\
\texttt{.dim\_C(q)} & computed & dim C\textasciicircum{}q\_eq (available even in the memory-lean mode) \\
\texttt{.invariant\_factors(q)} & computed & H\textasciicircum{}q = (+)\_i Z/ell\textasciicircum{}\{e\_i\}: the exponents e\_i, descending (invariant\_factors\_H\textasciicircum{}q) \\
\texttt{.r\_K(q)} & computed & r\_K(d\textasciicircum{}q) for q in \{0, 1\}: the K-mode's rank of the coboundary \\
\texttt{.ladic\_type(q)} & computed & \{free\_rank, torsion\} of H\textasciicircum{}q(C)\_(ell); defined only at a certified level (else the object's named std::logic\_error) \\
\texttt{.red(q)} & computed & red\_H\textasciicircum{}q: H\textasciicircum{}q(Z/ell\textasciicircum{}k) -\textgreater{} H\textasciicircum{}q(F\_ell), the (c5) matrix from the cyclic summands to the F\_ell object's D5b basis, over F\_ell \\
\end{membertable}
\octype{HeckeOperator}{T\_p (or U\_ell, the same path) on H\textasciicircum{}*: the coset tables, the certified tau datum and the assembled Theta / A\textasciicircum{}\{(q)\}}{\index{HeckeOperator@\texttt{HeckeOperator} (type)}}
\begin{membertable}
\texttt{.p} & structural & the Hecke prime \\
\texttt{.r} & structural & the exponent of N = p\textasciicircum{}r \\
\texttt{.cosets} & computed & the number of cosets m \\
\texttt{.certificates} & computed & the in-process certificate bundle compute() already ran (O12(a) and the counters) \\
\texttt{.json} & computed & the operator's own report JSON (levels (a) + (b)) \\
\texttt{.dims} & computed & dim H\textasciicircum{}q of the underlying complex \\
\texttt{.coefficient\_ring} & structural & the ring's canonical text: F\_ell, or Z/ell\textasciicircum{}\textless{}k\textgreater{} ([E] sec:ring (a)) \\
\texttt{.charpoly\_free} & computed & over Z/ell\textasciicircum{}k: \{precision, coeffs\} -- the charpoly of A\textasciicircum{}\{(1)\}(T\_p) on the free quotient, over Z/ell\textasciicircum{}precision, coefficients low -\textgreater{} high; defined only at a certified level \\
\texttt{.A(q)} & computed & A\textasciicircum{}\{(q)\}(T\_p), the cohomology-level matrix in degree q (over Z/ell\textasciicircum{}k: the (c5) matrix on the cyclic summands, a List of rows, entry (i, j) in [0, ell\textasciicircum{}\{e\_i\})) \\
\texttt{.Theta(q)} & computed & Theta\textasciicircum{}\{(q)\}\_flat on the flat cochains C\textasciicircum{}q(K; V) \\
\texttt{.Theta\_eq(q)} & computed & Theta\textasciicircum{}\{(q)\} restricted to C\textasciicircum{}q\_eq in the stored C\_eq basis \\
\texttt{.apply(q,\allowbreak{} v)} & computed & Theta\textasciicircum{}\{(q)\} applied to a flat cochain (the fused route) \\
\texttt{.apply\_route2(q,\allowbreak{} v)} & computed & the INDEPENDENT second route of [C, eq:heckechain] on the same cochain (matrix-free; it does not read the fused tables) \\
\end{membertable}
\octype{Family}{the Hecke operator family on one cohomology: T\_p, T\_\{p\textasciicircum{}r\} (recursion and the independent direct route), T\_n and the diamond T(p,p)}{\index{Family@\texttt{Family} (type)}}
\begin{membertable}
\texttt{.T(p)} & computed & T\_p per degree (this IS U\_ell at p = ell) \\
\texttt{.T\_pow(p,\allowbreak{} r)} & computed & the FULL T\_\{p\textasciicircum{}r\} by the algebra recursion of [C, thm:relations] \\
\texttt{.T\_pow\_direct(p,\allowbreak{} r)} & computed & the same operator by the INDEPENDENT double-coset route -- what T\_pow is cross-checked against \\
\texttt{.T\_pow\_primitive(p,\allowbreak{} r)} & computed & the PRIMITIVE double-coset operator T'\_\{p\textasciicircum{}r\} (equal to T\_p at r = 1; it differs from the standard T\_\{p\textasciicircum{}r\} for r \textgreater{}= 2) \\
\texttt{.T\_n(n)} & computed & T\_n for a composite index n \\
\texttt{.T\_pp(p)} & computed & the diamond T(p,p) = p\textasciicircum{}n I ([C, prop:diamond]) \\
\texttt{.theta(p,\allowbreak{} q)} & computed & Theta\textasciicircum{}\{(q)\}\_flat of the family's r = 1 pipeline for p \\
\end{membertable}
\octype{ContragredientLocalSystem}{D\_n = Sym\textasciicircum{}n(F\_ell\textasciicircum{}2)\textasciicircum{}vee on the 7-vertex model (rho\_n\textasciicircum{}vee(h) = rho\_n(adj h)\textasciicircum{}T), its cohomology, Phi\_n : M\_n -\textgreater{} D\_n and the Hecke action A(D\_n)(T\_p) (E7; lazy)}{\index{ContragredientLocalSystem@\texttt{ContragredientLocalSystem} (type)}}
\begin{membertable}
\texttt{.n} & structural & the symmetric power \\
\texttt{.p} & structural & the coefficient prime ell \\
\texttt{.d\_V} & structural & dim D\_n = n + 1 \\
\texttt{.coefficient\_ring} & structural & F\_ell (E7 is a field object) \\
\texttt{.phi\_is\_isomorphism} & structural & whether Phi\_n (x) F\_ell is an isomorphism, i.e. n \textless{} ell ([C] lem:symdual) \\
\texttt{.dims} & computed & dim H\textasciicircum{}q(Gamma; D\_n) for every degree q (golden dim\_H\_Dn\textasciicircum{}q) \\
\texttt{.dims\_M} & computed & dim H\textasciicircum{}q(Gamma; M\_n), the standard module, for comparison \\
\texttt{.Phi} & computed & Phi\_n : M\_n -\textgreater{} D\_n, (n+1) x (n+1) over F\_ell \\
\texttt{.induced\_Phi(q)} & computed & P\textasciicircum{}q = H\textasciicircum{}q(Phi\_n) : H\textasciicircum{}q(M\_n) -\textgreater{} H\textasciicircum{}q(D\_n) \\
\texttt{.A(p,\allowbreak{} q)} & computed & A(D\_n)(T\_p)\textasciicircum{}q, the Hecke action on H\textasciicircum{}q(D\_n) (U\_ell at p = ell); charpoly(D.A(p, q)) is the golden charpoly\_A\_Dn(T\_p)\textasciicircum{}q \\
\texttt{.A\_M(p,\allowbreak{} q)} & computed & A(M\_n)(T\_p)\textasciicircum{}q on the standard module, built from the SAME coset / tau data \\
\texttt{.intertwines(p)} & computed & \{q0, q1, q2\}: whether A(D\_n)(T\_p) P\textasciicircum{}q == P\textasciicircum{}q A(M\_n)(T\_p) exactly (UNTWISTED, factor 1; [C] prop:pairing) \\
\end{membertable}
\octype{Pairing}{the cone-level pairing H\textasciicircum{}t\_c(M\_n) x H\textasciicircum{}\{2-t\}(D\_n) -\textgreater{} F\_ell and the circle pairing: Grams, perfectness, the PD adjointness gate of [C] cl:pd (E7; lazy)}{\index{Pairing@\texttt{Pairing} (type)}}
\begin{membertable}
\texttt{.coefficient\_ring} & structural & F\_ell (E7 is a field object) \\
\texttt{.certificate} & computed & the pairing certificate re-derived: \{failing, failing\_items\} \\
\texttt{.perfect(t)} & computed & whether the cone pairing in H\_c degree t is perfect (golden pairing\_perfect\textasciicircum{}t) \\
\texttt{.gram(t)} & computed & G\textasciicircum{}t: dim H\textasciicircum{}t\_c x dim H\textasciicircum{}\{2-t\}(D\_n) (golden pairing\_matrix\textasciicircum{}t) \\
\texttt{.boundary\_perfect(s)} & computed & whether the circle pairing in degree s is perfect \\
\texttt{.boundary\_gram(s)} & computed & the circle Gram in degree s \\
\texttt{.pd\_gate(primes)} & computed & \{t0, t1, t2\}: A\_c\textasciicircum{}\{(t)T\} G\textasciicircum{}t == G\textasciicircum{}t A(D\_n)\textasciicircum{}\{(2-t)\} for EVERY p in \texttt{primes} (factor 1; t = 0 vacuous and true); at primes [2, 3, 5] it is the golden pd\_adjoint\textasciicircum{}t. A TEST of [C] cl:pd, which is open \\
\end{membertable}
\octype{Block}{one primary block of a decomposition: its dimension, residue degree, the per-prime minimal / characteristic polynomials and generalized-Jordan partition, and the E0b locality certificate}{\index{Block@\texttt{Block} (type)}}
\begin{membertable}
\texttt{.dim} & computed & dim\_F of the block \\
\texttt{.residue\_degree} & computed & f = lcm\_j deg(pi\_j) \\
\texttt{.certified} & computed & whether f is CERTIFIED (a witness proves it), not merely a lower bound \\
\texttt{.local} & computed & the E0b locality certificate: this block is ONE maximal ideal \\
\texttt{.phase} & computed & which phase reached the verdict (linear / berlekamp / residue\_search) \\
\texttt{.w\_hat\_exponent} & computed & N: the Teichmueller exponent of the witness \\
\texttt{.witness\_poly} & computed & pi\_w, the certifying witness's minimal polynomial \\
\texttt{.basis} & computed & the d x dim matrix whose columns span the block \\
\texttt{.primes} & computed & the family's primes, ascending \\
\texttt{.witness} & computed & the linear residue witness as [[p, c], ...] (empty in strict mode, where the witness is a matrix) \\
\texttt{.min\_poly(p)} & computed & the minimal polynomial of T\_p on this block \\
\texttt{.char\_poly(p)} & computed & the characteristic polynomial of T\_p on this block \\
\texttt{.character(p)} & computed & P\_p, the character polynomial of the locality certificate \\
\texttt{.jordan(p)} & computed & the generalized-Jordan partition of T\_p on this block, parts descending \\
\end{membertable}
\octype{Decomposition}{the certified primary decomposition of H\textasciicircum{}q as a module over the commutative F\_ell-algebra generated by a commuting family of Hecke matrices}{\index{Decomposition@\texttt{Decomposition} (type)}}
\begin{membertable}
\texttt{.q} & structural & the cohomological degree \\
\texttt{.ell} & structural & the coefficient prime \\
\texttt{.n} & structural & n\_sym of the coefficient system \\
\texttt{.primes} & structural & the family's primes, in the order given \\
\texttt{.blocks} & computed & the primary blocks, in canonical order \\
\texttt{.num\_blocks} & computed & how many primary blocks there are \\
\texttt{.U} & computed & the change of basis: the block bases, in block order \\
\texttt{.witness\_hash} & computed & the golden \texttt{decomp\_*\_witness} hash (the residue-witness data), as 0x... text \\
\texttt{.locality} & computed & the E0b locality loop's state: mode, budget, splits, completeness \\
\texttt{.json} & computed & the decomposition's own JSON record \\
\texttt{.certify()} & computed & the UNCONDITIONAL certificate: 0 iff every claim re-derives (never trusting the randomized factor search) \\
\texttt{.conjugated(p)} & computed & U\textasciicircum{}\{-1\} A\textasciicircum{}\{(q)\}(T\_p) U, verified exactly block-diagonal \\
\end{membertable}
\octype{Report}{the hecke\_report/v2 record of one degree: level (b) A\textasciicircum{}\{(q)\}(T\_p), level (c) the decomposition, and optionally level (a) Theta\textasciicircum{}\{(q)\}\_flat}{\index{Report@\texttt{Report} (type)}}
\begin{membertable}
\texttt{.schema} & structural & the frozen schema name \\
\texttt{.level\_a} & structural & whether level (a) was included \\
\texttt{.path} & structural & the file it was written to, or \textquotedbl{}\textquotedbl{} \\
\texttt{.degree} & computed & the degree q this record is about \\
\texttt{.q\_max} & computed & the top degree of THIS complex \\
\texttt{.truncated} & computed & false on a finite model \\
\texttt{.q\_max\_source} & computed & \textquotedbl{}derived\textquotedbl{} unless a setting pinned q\_max \\
\texttt{.degrees} & computed & [0 .. q\_max] \\
\texttt{.json} & computed & the report text \\
\texttt{.write(path)} & computed & write the report to a file (an empty path writes to stdout) \\
\texttt{.has\_key(key)} & computed & whether the record carries a given JSON key -- the schema assertion a transcript can make without freezing the whole text \\
\end{membertable}
\octype{CampaignSpec}{a parsed campaign spec (docs/matrix-engine/campaign.md) -- campaign\_spec(path), campaign\_spec\_text(text)}{\index{CampaignSpec@\texttt{CampaignSpec} (type)}}
\begin{membertable}
\texttt{.name} & structural & the spec's \texttt{name} \\
\texttt{.output\_dir} & structural & where run\_campaign writes by default \\
\texttt{.jobs} & structural & the number of jobs the plan has \\
\end{membertable}
\octype{Sturm}{the computed Sturm bound B'(k, ell) of [E] def:bprime: the least prime bound at which the F\_ell-ALGEBRA generated by T\_q, q \textless{}= B, reaches dim M\_k}{\index{Sturm@\texttt{Sturm} (type)}}
\begin{membertable}
\texttt{.k} & structural & the weight k = n + 2 \\
\texttt{.ell} & structural & the coefficient prime \\
\texttt{.outcome} & computed & found / none\_at\_cap / in\_progress \\
\texttt{.found} & computed & whether B' exists at or below the cap \\
\texttt{.B\_prime} & computed & B', meaningful iff \texttt{found} \\
\texttt{.cap} & computed & the cap in force (12*B(k) by default) \\
\texttt{.dim\_M\_k} & computed & dim M\_k(F\_ell) \\
\texttt{.closure\_dim} & computed & the dimension the algebra closure reached \\
\texttt{.primes\_done} & computed & the primes processed, ascending \\
\texttt{.certified} & computed & SturmBoundCertificate::verify() == 0, re-derived from the reported data \\
\texttt{.verify} & computed & the certificate's code (0 = every clause re-derives) \\
\texttt{.json} & computed & the certificate's JSON, \texttt{resumable} object included \\
\end{membertable}
\octype{Atlas}{Sigma(ell): the anemic level-1 eigensystems over weights \textless{}= ell + 1, their theta-twists, and the separating set S\_sep(ell) (HeckeAtlas, in process at ell in \{5, 7, 11, 13\}); lazy}{\index{Atlas@\texttt{Atlas} (type)}}
\begin{membertable}
\texttt{.ell} & structural & the prime \\
\texttt{.n\_max} & computed & the largest even weight cell n \textless{}= ell - 1 \\
\texttt{.primes} & computed & the anemic pool T \\
\texttt{.num\_systems} & computed & \textbar{}Sigma(ell)\textbar{} \\
\texttt{.chi\_count} & computed & distinct chi\_\{a,b\} as T-systems \\
\texttt{.chi\_in\_sigma} & computed & of which in Sigma \\
\texttt{.separating\_primes} & computed & the greedy S\_sep(ell) \\
\texttt{.separates} & computed & S\_sep separates Sigma u chi \\
\texttt{.s0\_separates} & computed & \{2,3,5,7\} \textbackslash{} \{ell\} already separates \\
\texttt{.cells} & computed & the weight cells \{n, dims, ideals, local\_ideals\} \\
\texttt{.json} & computed & the hecke\_atlas/v1 store text \\
\texttt{.system(i)} & computed & the i-th system \{id, r, type, base, twist, chi, eigenvalues, collected, in\_H1\} \\
\texttt{.lookup(eigenvalues)} & computed & the id of the system with these eigenvalues on \texttt{primes} (Ints at r = 1, coefficient lists at r \textgreater{} 1; any Frobenius conjugate), or -1 \\
\end{membertable}
\octype{GF}{the project-canonical field F\_ell[Y]/(C\_\{ell,r\}) ([E] def:canon; HeckeCanonicalField), C chosen by the \texttt{field\_polynomial} setting}{\index{GF@\texttt{GF} (type)}}
\begin{membertable}
\texttt{.ell} & structural & the prime \\
\texttt{.degree} & structural & r = [F : F\_ell] \\
\texttt{.polynomial} & structural & C\_\{ell,r\} \\
\texttt{.canonical} & structural & C is the CanonicalLess-least monic irreducible \\
\texttt{.zero} & structural & 0 \\
\texttt{.one} & structural & 1 \\
\texttt{.gen} & structural & Y, the class of the variable \\
\texttt{.elt(coeffs)} & structural & the element sum c\_i Y\textasciicircum{}i (coefficients low -\textgreater{} high, reduced mod C) \\
\texttt{.element\_at(i)} & structural & the i-th element of the field's fixed enumeration \\
\end{membertable}
\octype{GFElt}{an element of a GF field (coefficient vector in the basis 1, Y, ..., Y\textasciicircum{}\{r-1\})}{\index{GFElt@\texttt{GFElt} (type)}}
\begin{membertable}
\texttt{.coeffs} & structural & the coefficients, low -\textgreater{} high, width r \\
\texttt{.is\_zero} & structural & x == 0 \\
\texttt{.field} & structural & the GF handle \\
\texttt{.poly} & structural & the representative polynomial in Y \\
\end{membertable}
\octype{Settings}{a snapshot of the MATHEMATICAL settings (hecke::MathSettings, [X] Sec. 1.5): every registered key with its value and source, and the fingerprint}{\index{Settings@\texttt{Settings} (type)}}
\begin{membertable}
\texttt{.fingerprint} & structural & FNV-1a-64 of the non-default entries (0x0 at the defaults) \\
\texttt{.is\_default} & structural & every key at its default \\
\texttt{.keys} & structural & the registered keys, in registry order \\
\texttt{.values} & structural & \{key: value text\} \\
\texttt{.text} & structural & the canonical text \\
\texttt{.json} & structural & the settings JSON (values and sources) \\
\texttt{.get(key)} & structural & the value text of one key \\
\end{membertable}
\octype{ExtRecord}{one cell's hecke\_ext/v5 record (HeckeExtDriver, Stages 1-5): its maximal ideals, groups, u\_ell verdict and settings; lazy}{\index{ExtRecord@\texttt{ExtRecord} (type)}}
\begin{membertable}
\texttt{.n} & structural & the symmetric power n \\
\texttt{.ell} & structural & the coefficient prime \\
\texttt{.schema} & computed & the record schema \\
\texttt{.weight} & computed & k = n + 2 \\
\texttt{.coefficient\_ring} & computed & the record's ring (F\_ell) \\
\texttt{.dims} & computed & dim H\textasciicircum{}q of the bulk, every degree \\
\texttt{.degrees\_run} & computed & the bulk degrees Stages 2-5 ran on \\
\texttt{.operators} & computed & the primes whose operators Stage 1 computed \\
\texttt{.anemic\_pool} & computed & the anemic pool \{q \textless{}= 31\} \textbackslash{} \{ell\} \\
\texttt{.s\_sep} & computed & S\_sep(ell) as Stage 2(3) ran it (empty if not computed) \\
\texttt{.dim\_M\_k} & computed & dim M\_k \\
\texttt{.dim\_S\_k} & computed & dim S\_k \\
\texttt{.dim\_dickson} & computed & the Dickson count \\
\texttt{.H1\_shriek} & computed & dim H\textasciicircum{}1\_! (cone = true), else nil \\
\texttt{.les\_exact} & computed & the cone's LES exact (cone = true), else nil \\
\texttt{.groups} & computed & the record's groups, \{name, available, dim\} \\
\texttt{.u\_ell} & computed & the cell's U\_ell verdict and its resumable state \\
\texttt{.ideals} & computed & the maximal ideals, as MaximalIdeal handles \\
\texttt{.num\_ideals} & computed & the number of maximal ideals \\
\texttt{.has\_sign\_records} & computed & whether sign\_split = true built H\textasciicircum{}+ / H\textasciicircum{}- \\
\texttt{.ideals\_plus} & computed & the ideals of the H\textasciicircum{}+ record (sign\_split = true) \\
\texttt{.ideals\_minus} & computed & the ideals of the H\textasciicircum{}- record (sign\_split = true) \\
\texttt{.settings} & computed & the MathSettings the record was computed under \\
\texttt{.verify} & computed & HeckeExtRecord::verify() (0 = consistent) \\
\texttt{.json} & computed & the hecke\_ext/v5 record text (carries the build's commit) \\
\texttt{.digest} & computed & FNV-1a-64 of the record text with \texttt{commit} redacted (NOT the golden ext\_json\_* hash) \\
\texttt{.ideal(i)} & computed & the i-th maximal ideal \\
\end{membertable}
\octype{MaximalIdeal}{one maximal ideal m of a cell's Hecke algebra (an ideal of the ExtRecord)}{\index{MaximalIdeal@\texttt{MaximalIdeal} (type)}}
\begin{membertable}
\texttt{.id} & computed & the ideal's id in its record \\
\texttt{.sign} & computed & all / plus / minus \\
\texttt{.dim} & computed & dim\_F\_ell of the block H\_m \\
\texttt{.group} & computed & the ambient group of the block \\
\texttt{.local} & computed & certified local \\
\texttt{.residue\_degree} & computed & f = [T/m : F\_ell] \\
\texttt{.residue\_degree\_certified} & computed & whether f is certified \\
\texttt{.witness\_poly} & computed & pi\_w, the witness's minimal polynomial \\
\texttt{.w\_hat\_exponent} & computed & N of the Teichmueller lift (local only) \\
\texttt{.final} & computed & true / false / heuristic \\
\texttt{.final\_reason} & computed & why not final ('' when final) \\
\texttt{.certificate} & computed & cyclic / commutant / sturm(B') / dickson / boundary / unavailable \\
\texttt{.certificate\_routes} & computed & the routes' outcomes \{a, b, c, d0, d1\} (d0 / d1: the finality certificates, hecke\_ext/v4) \\
\texttt{.generators\_certified} & computed & the certified generating primes \\
\texttt{.boundary\_reps} & computed & R(ell), the class representatives (d1) tested (empty unless it reached them) \\
\texttt{.operator\_primes} & computed & the primes of the restricted operators \\
\texttt{.tangent\_active} & computed & the primes whose n\_p is active in m/m\textasciicircum{}2 \\
\texttt{.embdim} & computed & dim\_K m/m\textasciicircum{}2 (local only) \\
\texttt{.nilpotency} & computed & the nilpotency index nu (local only) \\
\texttt{.gorenstein} & computed & dim soc = f (local only) \\
\texttt{.label} & computed & the atlas label \{type, system, theta\_twist, ...\} or nil \\
\texttt{.predicted} & computed & the predictions attached to the ideal \\
\texttt{.unpredicted} & computed & what was found and not predicted \\
\texttt{.modules} & computed & the LoewyProfile of every group the ideal lives in \\
\texttt{.algebra} & computed & the local Hecke algebra T\_m (HeckeAlgebra) \\
\texttt{.t(p)} & computed & the restricted operator t\_p = T\_p\textbar{}\_\{H\_m\} \\
\texttt{.character(p)} & computed & P\_p, the eigenvalue of T\_p as a polynomial in the witness (local only) \\
\end{membertable}
\octype{HeckeAlgebra}{the local Hecke algebra T\_m = A\_S of a certified-local ideal (HeckeLocalAlgebra; [X] Sec. 7's \texttt{LocalAlgebra})}{\index{HeckeAlgebra@\texttt{HeckeAlgebra} (type)}}
\begin{membertable}
\texttt{.ell} & computed & the prime \\
\texttt{.D} & computed & dim\_F\_ell of the block \\
\texttt{.f} & computed & the residue degree \\
\texttt{.dim} & computed & dim\_F\_ell A\_S \\
\texttt{.dim\_m} & computed & dim\_F\_ell m\_S \\
\texttt{.nilpotency} & computed & nu: m\textasciicircum{}nu = 0 != m\textasciicircum{}\{nu-1\} \\
\texttt{.hilbert} & computed & dim(m\textasciicircum{}i/m\textasciicircum{}\{i+1\})/f \\
\texttt{.socle\_dim} & computed & dim\_F\_ell soc A\_S \\
\texttt{.gorenstein} & computed & dim soc A\_S == f \\
\texttt{.embdim} & computed & dim\_K m/m\textasciicircum{}2 \\
\texttt{.primes} & computed & S, ascending \\
\texttt{.active\_set} & computed & the primes whose n\_p-bar is nonzero \\
\texttt{.minimal\_active} & computed & a K-spanning subset, greedy in ascending p \\
\texttt{.min\_generators} & computed & dim(V/m V)/f \\
\texttt{.cyclic} & computed & certificate (a): V cyclic over A\_S \\
\texttt{.verify} & computed & certificate (a)'s verify() (0 = re-derives) \\
\texttt{.basis} & computed & an F\_ell-basis of A\_S (canonical RREF) \\
\texttt{.radical} & computed & an F\_ell-basis of m\_S \\
\texttt{.w\_hat} & computed & the Teichmueller lift of the witness \\
\texttt{.invariants} & computed & the SAME invariants by the certificate-free route of algebra\_invariants(t\_p) -- a cross-check \\
\texttt{.json} & computed & the LocalAlgebra JSON \\
\texttt{.t(p)} & computed & t\_p = T\_p restricted to the block \\
\texttt{.n(p)} & computed & n\_p = t\_p - P\_p(w\_hat), the nilpotent generator \\
\end{membertable}
\octype{LoewyProfile}{the Stage-5 module structure of a block at m (HeckeModuleStructure; [X] Sec. 7's \texttt{ModuleStructure}): Loewy / socle series, Jordan types, structure maps}{\index{LoewyProfile@\texttt{LoewyProfile} (type)}}
\begin{membertable}
\texttt{.group} & structural & the ambient group \\
\texttt{.D} & computed & dim\_F\_ell V \\
\texttt{.f} & computed & the residue degree \\
\texttt{.loewy} & computed & Loewy layer dimensions in units of f \\
\texttt{.socle} & computed & socle layer dimensions in units of f \\
\texttt{.m\_dims} & computed & dim\_F\_ell m\textasciicircum{}j V, j = 0 .. nu\_V \\
\texttt{.socle\_dims} & computed & dim\_F\_ell soc\textasciicircum{}j V \\
\texttt{.primes} & computed & the generators' primes \\
\texttt{.cyclic} & computed & V cyclic over T\_m \\
\texttt{.free} & computed & V free of rank one \\
\texttt{.min\_generators} & computed & dim(V/mV)/f \\
\texttt{.semisimple} & computed & mV = 0 (no extension at m) \\
\texttt{.verify} & computed & ModuleStructure::verify() (0 = re-derives) \\
\texttt{.json} & computed & the module structure JSON \\
\texttt{.structure\_maps} & computed & the ExtensionClass of every consecutive layer pair \\
\texttt{.jordan(p)} & computed & the Jordan type of n\_p on V, in units of f \\
\texttt{.layer(j)} & computed & the canonical basis of the Loewy layer m\textasciicircum{}j V / m\textasciicircum{}\{j+1\} V (columns) \\
\texttt{.structure\_map(j)} & computed & the ExtensionClass of the layer pair (j, j+1) \\
\end{membertable}
\octype{ExtensionClass}{the structure map theta\_j between two consecutive Loewy layers ([E] D4): the extension class the generators n\_p glue the layers by}{\index{ExtensionClass@\texttt{ExtensionClass} (type)}}
\begin{membertable}
\texttt{.top\_layer} & computed & j \\
\texttt{.bottom\_layer} & computed & j + 1 \\
\texttt{.top\_dim} & computed & the top layer's K-dimension \\
\texttt{.bottom\_dim} & computed & the bottom layer's K-dimension \\
\texttt{.rank} & computed & F\_ell rank of the stacked map \\
\texttt{.rank\_over\_f} & computed & its K-rank \\
\texttt{.active} & computed & the primes whose own map is nonzero \\
\texttt{.rank\_p} & computed & the F\_ell rank per generator \\
\texttt{.split} & computed & rank == 0: the two layers are not glued \\
\texttt{.matrix(p)} & computed & the matrix of n\_p between the layers (rows bottom, columns top) \\
\end{membertable}
\octype{Presentation}{generators and a minimal relation set of a certified-local block over T\_m (HeckePresentation, [X] Sec. 2.2 Stage 5, OQ11)}{\index{Presentation@\texttt{Presentation} (type)}}
\begin{membertable}
\texttt{.available} & structural & whether the block was local \\
\texttt{.min\_generators} & structural & t, over A\_S \\
\texttt{.min\_relations} & structural & over A\_S \\
\texttt{.dim\_relation\_module} & structural & t dim A\_S - D \\
\texttt{.free} & structural & no relations \\
\texttt{.generator\_indices} & structural & the generators' standard indices \\
\texttt{.generators} & structural & the generators as columns \\
\texttt{.relations} & structural & the relations as t-tuples over the A\_S basis \\
\texttt{.verify} & structural & Presentation::verify() (0 = re-derives) \\
\texttt{.json} & structural & the presentation JSON \\
\end{membertable}
\octype{Involution}{the production involution iota\textasciicircum{}* = -iota\textasciicircum{}*\_geo on H ([C] def:iota(iv); HeckeInvolution) with the +/- split and, per degree, the sign split of the Hecke action and the Eisenstein line (HeckeSignSplit); lazy}{\index{Involution@\texttt{Involution} (type)}}
\begin{membertable}
\texttt{.n} & computed & the symmetric power \\
\texttt{.ell} & computed & the prime \\
\texttt{.status} & computed & the split's status \\
\texttt{.coefficient\_ring} & computed & F\_ell, or Z/ell\textasciicircum{}k (slot 20: iota honours the ring) \\
\texttt{.lengths\_plus} & computed & length H\textasciicircum{}q\_+ per degree (= dims\_plus over F\_ell) \\
\texttt{.lengths\_minus} & computed & length H\textasciicircum{}q\_- per degree (= dims\_minus over F\_ell) \\
\texttt{.dims} & computed & dim H\textasciicircum{}q \\
\texttt{.dims\_plus} & computed & dim H\textasciicircum{}q\_+ \\
\texttt{.dims\_minus} & computed & dim H\textasciicircum{}q\_- \\
\texttt{.certificate} & computed & InvolutionCertificate::verify() (0 = obl:iota discharged; over Z/ell\textasciicircum{}k: iota\textasciicircum{}2 = id and H = H\textasciicircum{}+ (+) H\textasciicircum{}- were ENGINE\_CHECKed) \\
\texttt{.json} & computed & hecke\_involution/v1 (over Z/ell\textasciicircum{}k: hecke\_involution\_ring/v1) \\
\texttt{.invariant\_factors\_plus(q)} & computed & H\textasciicircum{}q\_+ = (+)\_i Z/ell\textasciicircum{}\{e\_i\}: the e\_i, descending (F\_ell: all 1) \\
\texttt{.invariant\_factors\_minus(q)} & computed & H\textasciicircum{}q\_- = (+)\_i Z/ell\textasciicircum{}\{e\_i\}: the e\_i, descending (F\_ell: all 1) \\
\texttt{.matrix(q)} & computed & iota\textasciicircum{}* on H\textasciicircum{}q \\
\texttt{.plus(q)} & computed & a basis of H\textasciicircum{}q\_+ (columns, in H\textasciicircum{}q coordinates) \\
\texttt{.minus(q)} & computed & a basis of H\textasciicircum{}q\_- (columns, in H\textasciicircum{}q coordinates) \\
\texttt{.sign\_split(q)} & computed & the Hecke action's sign split at \{2,3,5,7\} \textbackslash{} \{ell\} and the Eisenstein line \\
\texttt{.eisenstein\_line(q)} & computed & the Eisenstein line in the stored H\textasciicircum{}q basis (defined only) \\
\end{membertable}
\octype{BoundaryCohomology}{H\textasciicircum{}*\_bd on the boundary circle (SL2ZBoundaryPair), with the restriction r and T\_p\textasciicircum{}bd (HeckeBoundaryOperator); lazy}{\index{BoundaryCohomology@\texttt{BoundaryCohomology} (type)}}
\begin{membertable}
\texttt{.n} & computed & the symmetric power \\
\texttt{.ell} & computed & the prime \\
\texttt{.model} & computed & the boundary model \\
\texttt{.is\_group\_cohomology} & computed & 2 in R\textasciicircum{}x (the circle computes group cohomology) \\
\texttt{.dims} & computed & dim H\textasciicircum{}q\_bd \\
\texttt{.bulk\_dims} & computed & dim H\textasciicircum{}q of the bulk \\
\texttt{.q\_max} & computed & the top degree of THIS complex (derived) \\
\texttt{.truncated} & computed & false on a finite model \\
\texttt{.q\_max\_source} & computed & 'derived' unless truncated (E9) \\
\texttt{.degrees} & computed & [0 .. q\_max], derived \\
\texttt{.coefficient\_ring} & computed & the complex's ring (the extension layer is F\_ell) \\
\texttt{.r(q)} & computed & the restriction H\textasciicircum{}q -\textgreater{} H\textasciicircum{}q\_bd \\
\texttt{.A(p,\allowbreak{} q)} & computed & T\_p\textasciicircum{}bd on H\textasciicircum{}q\_bd \\
\end{membertable}
\octype{ConeCohomology}{H\textasciicircum{}*\_c = H\textasciicircum{}*(Cone(r)) with its long exact sequence (EquivariantMappingCone) and, per prime, the cone action A\_c(T\_p) with h DERIVED from the certified pair homotopy eta over Z (HeckeConeAction); lazy}{\index{ConeCohomology@\texttt{ConeCohomology} (type)}}
\begin{membertable}
\texttt{.n} & computed & the symmetric power \\
\texttt{.ell} & computed & the prime \\
\texttt{.name} & computed & the cone's name \\
\texttt{.dims} & computed & dim H\textasciicircum{}q\_c \\
\texttt{.bulk\_dims} & computed & dim H\textasciicircum{}q of the bulk \\
\texttt{.boundary\_dims} & computed & dim H\textasciicircum{}q\_bd \\
\texttt{.q\_max} & computed & the top degree of THIS complex (derived) \\
\texttt{.truncated} & computed & false on a finite model \\
\texttt{.q\_max\_source} & computed & 'derived' unless truncated (E9) \\
\texttt{.degrees} & computed & [0 .. q\_max], derived \\
\texttt{.coefficient\_ring} & computed & the complex's ring (the extension layer is F\_ell) \\
\texttt{.certificate} & computed & ConeCertificate::verify() (0 = re-derives) \\
\texttt{.H1\_shriek} & computed & dim H\textasciicircum{}1\_! = dim ker r\textasciicircum{}1 (nil if not derivable) \\
\texttt{.les} & computed & the long exact sequence spots \{term, q, dim, rank\_in, rank\_out, asserted, exact\} \\
\texttt{.les\_exact} & computed & every asserted spot of the LES is exact \\
\texttt{.pi(q)} & computed & H\textasciicircum{}q\_c -\textgreater{} H\textasciicircum{}q \\
\texttt{.r(q)} & computed & H\textasciicircum{}q -\textgreater{} H\textasciicircum{}q\_bd \\
\texttt{.iota(q)} & computed & H\textasciicircum{}q\_bd -\textgreater{} H\textasciicircum{}\{q+1\}\_c (the connecting map) \\
\texttt{.A(p,\allowbreak{} q)} & computed & A\_c(T\_p) on H\textasciicircum{}q\_c \\
\texttt{.h(p,\allowbreak{} q)} & computed & the cochain homotopy h at p in degree q, DERIVED from eta (not stored) \\
\texttt{.pair\_homotopy(p)} & computed & the certified pair homotopy eta at p: \{eta\_zero, slots, certify, solve\_regime, ...\} \\
\texttt{.action\_certificate(p)} & computed & ConeActionCertificate::verify() of the action at p (0 = re-derives) \\
\end{membertable}
\octype{DicksonLES}{the long exact sequence of 0 -\textgreater{} Sym\textasciicircum{}n (x) det -\textgreater{} Sym\textasciicircum{}\{n+ell+1\} -\textgreater{} Q\_n -\textgreater{} 0 on the model over F\_ell, ell \textgreater{}= 5 ([E] sec:thetaobj; HeckeDicksonTheta): the six terms, the five maps in the stored bases (theta = L\_* source-indexed, delta by the snake construction), exactness at all six spots as subspaces, H\textasciicircum{}0 the Dickson ring; lazy}{\index{DicksonLES@\texttt{DicksonLES} (type)}}
\begin{membertable}
\texttt{.n} & structural & the symmetric power of the source term \\
\texttt{.ell} & structural & the prime \\
\texttt{.n\_up} & structural & n + ell + 1, the middle term's symmetric power \\
\texttt{.dims} & computed & \{H0\_n, H0\_up, H0\_Q, H1\_n, H1\_up, H1\_Q\} \\
\texttt{.ranks} & computed & \{theta0, pi0, delta, theta1, pi1\} \\
\texttt{.maps} & computed & \{theta0, pi0, delta, theta1, pi1\} as Matrix values in the stored bases (column convention) \\
\texttt{.exact} & computed & exactness certified at all six spots as subspaces \\
\texttt{.ker\_theta} & computed & a basis of ker theta\textasciicircum{}1 (columns, in H1\_n's stored coordinates) \\
\texttt{.im\_delta} & computed & a basis of im delta (columns, in H1\_n's stored coordinates) \\
\texttt{.ker\_equals\_im} & computed & ker theta\textasciicircum{}1 = im delta as subspaces of H1\_n \\
\texttt{.dickson\_count} & structural & \#\{(a, b) : a(ell + 1) + b(ell\textasciicircum{}2 - ell) = n\} \\
\texttt{.certificate} & computed & DicksonThetaCertificate code (0 = every check of [E] (6a)-(6d) held) \\
\texttt{.coefficient\_ring} & structural & the complex's ring (theta is F\_ell only) \\
\end{membertable}
\octype{DescentRecord}{the integral descent record integral\_descent\_record/v1 of H\textasciicircum{}q(Gamma\_0(N); Sym\textasciicircum{}n Z\textasciicircum{}2) over Z ([C] thm:descentcert; slot-24 D6): the groups to q\_max = 3 (object = group, truncated, periodic from q = 2), the bad primes, the certified precision K per prime, char\_Z(T\_p \textbar{} H\textasciicircum{}1/tors), the torsion Hecke modules by their invariants, the boundary pair (sec:descentpair); only under coefficient\_ring = Z; lazy}{\index{DescentRecord@\texttt{DescentRecord} (type)}}
\begin{membertable}
\texttt{.N} & structural & the level \\
\texttt{.n} & structural & the symmetric power \\
\texttt{.primes} & structural & the Hecke primes P (none dividing N) \\
\texttt{.coefficient\_ring} & structural & Z (the record exists only under coefficient\_ring = Z) \\
\texttt{.object} & structural & group (group cohomology, never Bredon; OD-10 (a)) \\
\texttt{.q\_max} & structural & 3, derived from thm:descentcert(6) \\
\texttt{.truncated} & structural & true: degrees above q\_max are not reported (they are periodic, thm:descentcert(6)) \\
\texttt{.q\_max\_source} & structural & derived \\
\texttt{.degrees} & structural & [0, 1, 2, 3] \\
\texttt{.groups} & computed & H\textasciicircum{}0 .. H\textasciicircum{}3 as abelian groups (\textquotedbl{}Z\textasciicircum{}3 + Z/2 + Z/4\textquotedbl{}); empty when refused \\
\texttt{.bad\_primes} & computed & B(M\_N) of thm:descentcert(1): List of [ell, count] (prop:badprimesN(1)), ell \textgreater{}= 5 \\
\texttt{.torsion\_primes} & computed & the primes whose torsion is read: 2, 3 (n even) and B(M\_N) \\
\texttt{.K} & computed & per prime of the groups' run: \{ell, route, K, W, certified\} (certification provenance, not hashed) \\
\texttt{.charpoly} & computed & per p: \{p, r, coefficients (high -\textgreater{} low, String: an integer, \textquotedbl{}[sign,D digits]\textquotedbl{} beyond 2\textasciicircum{}63, \textquotedbl{}?\textquotedbl{} undetermined), exact (high -\textgreater{} low, String: every coefficient as an exact integer, \textquotedbl{}?\textquotedbl{} undetermined), determined, digest, refused\} -- char\_Z(T\_p \textbar{} H\textasciicircum{}1/tors), prop:crtcharpolyN's output rule \\
\texttt{.charpoly\_factored} & computed & per p: \{p, status (ok \textbar{} unavailable), expanded, factored (\textquotedbl{}(X + 24)\textasciicircum{}2 (X - 2049)\textquotedbl{}), factors: [\{coefficients (high -\textgreater{} low, String), degree, multiplicity, irreducible\}], complete, verified, note\} -- char\_Z(T\_p \textbar{} H\textasciicircum{}1/tors) factored over Z (Yun + Zassenhaus, exact; the product of the factors is checked against char\_Z); presentation only, not in the record's JSON or digest \\
\texttt{.torsion} & computed & per (q, ell): \{q, ell, exponents, module\_certified, ops: [\{p, invariants\}], refused\} -- the torsion Hecke module by its choice-free invariants (prop:torsionhecke(1)) \\
\texttt{.pair} & computed & the boundary pair (sec:descentpair): \{bulk, boundary, compact\} each \{q\_max, truncated, q\_max\_source, degrees, groups\}, q = 0 .. 4 (compact, the cone, also derived\_from = [bulk, boundary]); refused \\
\texttt{.groups\_inv} & computed & H\textasciicircum{}0 .. H\textasciicircum{}3 as structures: per q \{free\_rank, torsion: [[p, k, mult], ...]\} = Z\textasciicircum{}free\_rank + (+) (Z/p\textasciicircum{}k)\textasciicircum{}mult (the same groups as \texttt{groups}); empty when refused \\
\texttt{.pair\_inv} & computed & the boundary pair's groups as structures: \{bulk, boundary, compact\}, each a List over q = 0 .. 4 of \{free\_rank, torsion: [[p, k, mult], ...]\} (the same groups as \texttt{pair}) \\
\texttt{.status} & computed & ok \textbar{} refused \\
\texttt{.refusal} & computed & the first refused part's reason (\textquotedbl{}\textquotedbl{} when ok) \\
\texttt{.digest} & computed & record\_fnv1a64: FNV-1a-64 of the record's invariants and the settings' canonical text \\
\texttt{.json} & computed & the record as one JSON line (schema integral\_descent\_record/v1) \\
\texttt{.module(ell,\allowbreak{} K)} & structural & H\textasciicircum{}1(Gamma\_0(N); Sym\textasciicircum{}n Z\textasciicircum{}2) (x) Z/ell\textasciicircum{}K with T\_p (p in primes) by cochain-level lifts (slot-25 EXPORT; [C] prop:torsionhecke, prop:strictliftN): \{N, n, ell, K, K\_max, route, carrier, status (ok \textbar{} blocked(export) \textbar{} refused), refusal, free\_rank, torsion (exponents at K), invariant\_factors, generators: [\{kind, exponent, cocycle (a cocycle mod ell\textasciicircum{}K)\}], ops: [\{p, matrix (column convention, row i mod ell\textasciicircum{}exponent\_i), charpoly\_free (high -\textgreater{} low, mod ell\textasciicircum{}K)\}]\}; one run at the route's width serves every K \textless{}= K\_max (computed on the first call per ell; it does not force the record); ell in \{2, 3\}: the tree model, ell \textgreater{}= 5: the cokernel presentation (K\_max \textless{}= 1 for ell \textgreater{}= 46341) \\
\texttt{.module\_Z(mode)} & computed & H\textasciicircum{}1(Gamma\_0(N); Sym\textasciicircum{}n Z\textasciicircum{}2) over Z as a presented Z[T\_p]-module ([C] prop:torsionhecke): T\_p = [[tau, M], [0, F]] in generators torsion first, then free; module\_Z() = EXPORT-Z step 2 (F exact over Z), module\_Z(\textquotedbl{}local\textquotedbl{}) = step 1 (slot-25 EXPORT-Z, F pending). \{N, n, primes, status (ok \textbar{} blocked(export)), refusal, splitting (export\_retraction/v1), retractions: [\{ell, route, K\_run, K, fnv, [step 2: pi\_B (t\_ell x r, row a mod d\_a: the retraction on B, the torsion coordinates of the classes [b\_j]), pi\_B\_fnv]\}], torsion\_primes, largest\_torsion\_prime, invariant\_factors (the elementary divisors d\_i = ell\textasciicircum{}sigma\_i, ell ascending; key name kept), torsion\_generators: [\{ell, exponent, d\}], free\_rank (a blocked value prints these four as \textquotedbl{}blocked\textquotedbl{}), free\_block (step 2: \textquotedbl{}exact\textquotedbl{}, or \textquotedbl{}refused\textquotedbl{} with free\_refusal; local: \textquotedbl{}pending(step 2)\textquotedbl{}), M\_basis (step 2: global, M in the basis B; local: M's ell-block is against module(ell, K)'s free generators), entry\_encoding (step 2: exact, F as decimal integer strings; local: word), split (unknown), presentation, [step 2: basis: \{rule (iso\_dixon/v3), q0, saturation (det\_pm1), basis\_fnv, phi\_fnv, bits\_B, bits\_Phi, bound (phi0 \textbar{} phi\_cert by the Auto size rule; - at rank 0)\}, free\_refusal], ops: [\{p, tau (row i mod d\_i, block-diagonal over ell), M: [\{ell, K, matrix\}] (step 2: against the classes [b\_j] of B, T\_p b\_j = sum\_i F\_ij b\_i + sum\_a M\_aj g\_a, i.e. M = M\_local C - pi\_B F + tau pi\_B; local: against module(ell, K)'s free generators), F (step 2: r x r decimal strings; local: \textquotedbl{}pending(step 2)\textquotedbl{}), F\_local: [\{ell, modulus, matrix\}] (the export's free block mod ell\textasciicircum{}K, not F), charpoly\_Z (the record's char\_Z(T\_p \textbar{} H\textasciicircum{}1/tors), high -\textgreater{} low), [step 2: F\_fnv, F\_proof (apriori \textbar{} apost \textbar{} rank0), F\_primes, M\_fnv (FNV-1a-64 of M over every ell, span convention), C: [\{ell, modulus, matrix\}] (B in the export's local free generators: C F = F\_local C mod ell\textasciicircum{}K)]\}], torsion\_canon: [\{ell, ops: [\{p, invariants\}]\}] (the record's)\}; unique up to [[Aut T, Hom(L, T)], [0, GL\_r(Z)]] \\
\end{membertable}
